\item \textbf{Cálculo del momento de inercia de una sonda espacial.}
Ha obtenido un trabajo CE a tiempo parcial en una empresa que hace pequeñas sondas que son liberadas de los satélites para estudiar la atmósfera muy delgada en la ubicación de las órbitas de estos. Para mantener las sondas en una orientación adecuada en el espacio, se girarán alrededor de su eje antes de soltarse. Es importante conocer el momento de inercia de la sonda de forma extraña. Su jefe le pide que mida su momento de inercia. Configura un sistema como el de la figura 15.18, modificándolo añadiendo un marco muy liviano (figura P15.29) en el cual puede colocar objetos, centrándolos en el disco. El marco está unido a los bordes del disco. El disco es de masa $M = 5.25\text{ kg}$ y tiene un radio de $R = 25.8\text{ cm}$. Gira el disco vacío desde su posición de equilibrio y lo deja funcionar como un péndulo de torsión. Mide cuidadosamente su periodo de oscilación para ser $T_{\text{vacío}} = 10.8\text{ s}$.

A continuación, coloca la sonda en el disco y ajusta su posición hasta que el disco cuelgue exactamente en posición horizontal, para que sepa que el centro de masa de la sonda está directamente sobre el centro del disco. Gira el disco cargado desde su posición de equilibrio y lo deja funcionar como un péndulo de torsión.
\begin{enumerate}
    \item Mida cuidadosamente el periodo de oscilación para $T_{\text{cargado}} = 18.7\text{ s}$, y de este resultado determine el momento de inercia de la sonda sobre su centro de masa.
    \item Cuando presenta los resultados a su supervisora, ella le pregunta sobre el momento de inercia del marco que construyó. Vuelve a su escritorio y lo piensa. Cuando considera que el cuadro tiene algún momento de inercia, ¿el valor calculado en el inciso (1) es demasiado alto o demasiado bajo?
\end{enumerate}